\documentclass[12pt, a4paper]{article}
\usepackage[utf8]{inputenc}
\usepackage[T1]{fontenc}
\usepackage{amsmath, amssymb}
\usepackage{geometry}
\usepackage{booktabs}
\usepackage{titlesec}
\usepackage{xcolor}
\usepackage{mdframed}
\usepackage{enumitem}
\usepackage{array}
\usepackage{parskip}
\usepackage{listings}
\usepackage{courier}

\geometry{margin=2.5cm, top=3cm, bottom=3cm}

\definecolor{ballblue}{RGB}{0, 82, 147}
\definecolor{lightgray}{RGB}{245, 245, 245}
\definecolor{darkgray}{RGB}{80, 80, 80}
\definecolor{accentgray}{RGB}{180, 180, 180}
\definecolor{codegreen}{RGB}{40, 120, 40}
\definecolor{codegray}{RGB}{100, 100, 100}
\definecolor{codebg}{RGB}{250, 250, 250}
\definecolor{codeblue}{RGB}{0, 60, 180}
\definecolor{codeorange}{RGB}{180, 80, 0}

\titleformat{\section}
  {\large\bfseries\color{ballblue}}{}
  {0em}{}[\color{ballblue}\titlerule]
\titlespacing{\section}{0pt}{18pt}{8pt}

\titleformat{\subsection}
  {\normalsize\bfseries\color{darkgray}}{}
  {0em}{}
\titlespacing{\subsection}{0pt}{12pt}{4pt}

\mdfdefinestyle{destaque}{
  backgroundcolor=lightgray,
  linecolor=ballblue, linewidth=1.2pt,
  leftline=true, rightline=false, topline=false, bottomline=false,
  innerleftmargin=10pt, innerrightmargin=10pt,
  innertopmargin=8pt, innerbottommargin=8pt,
  skipabove=6pt, skipbelow=6pt,
  nobreak=true,
}

\mdfdefinestyle{codigo}{
  backgroundcolor=codebg,
  linecolor=accentgray, linewidth=0.8pt,
  leftline=true, rightline=false, topline=false, bottomline=false,
  linecolor=codegreen,
  innerleftmargin=10pt, innerrightmargin=10pt,
  innertopmargin=8pt, innerbottommargin=8pt,
  skipabove=6pt, skipbelow=6pt,
  nobreak=true,
}

\lstdefinelanguage{Julia}{
  keywords={using, function, for, in, if, end, return, true, false, @variable, @constraint, @objective, Model, optimize, Min},
  keywordstyle=\color{codeblue}\bfseries,
  comment=[l]{\#},
  commentstyle=\color{codegray}\itshape,
  stringstyle=\color{codeorange},
  morestring=[b]",
  basicstyle=\ttfamily\small,
  breaklines=true,
  showstringspaces=false,
}

\lstset{language=Julia}

% =============================================================================
\begin{document}

% ── Capa ─────────────────────────────────────────────────────────────────────
\begin{center}
  {\color{ballblue}\rule{\linewidth}{2pt}}\\[8pt]
  {\LARGE\bfseries Modelo de Programa\c{c}\~{a}o Linear}\\[4pt]
  {\large\color{darkgray} Planejamento Macrosc\'{o}pico --- Ball Corporation}\\[8pt]
  {\color{ballblue}\rule{\linewidth}{2pt}}
\end{center}

\vspace{10pt}
\begin{center}
  \small\color{darkgray}
  \textbf{1.~Minimiza\c{c}\~{a}o de custos}
  \enspace\textbullet\enspace
  \textbf{2.~Conjuntos}
  \enspace\textbullet\enspace
  \textbf{3.~Par\^{a}metros}
  \enspace\textbullet\enspace
  \textbf{4.~Fun\c{c}\~{a}o Auxiliar}
  \enspace\textbullet\enspace
  \textbf{5.~Vari\'{a}veis}
  \enspace\textbullet\enspace
  \textbf{6.~Fun\c{c}\~{a}o Objetivo}
  \enspace\textbullet\enspace
  \textbf{7.~Restri\c{c}\~{o}es}
  \enspace\textbullet\enspace
  \textbf{8.~Formulac\c{c}\~{a}o Compacta}
  \enspace\textbullet\enspace
  \textbf{9.~Implementa\c{c}\~{a}o}
  \enspace\textbullet\enspace
  \textbf{10.~An\'{a}lise de Sensibilidade}
\end{center}
\vspace{4pt}
{\color{accentgray}\rule{\linewidth}{0.5pt}}

% ============================================================
\section{Modelo de minimiza\c{c}\~{a}o de custos}

O modelo busca \textbf{minimizar o custo total de opera\c{c}\~{a}o} da malha industrial da Ball
Corporation, contemplando quatro componentes: produ\c{c}\~{a}o, estocagem, transfer\^{e}ncia
log\'{i}stica e penalidade por demanda n\~{a}o atendida. O modelo \'{e} estritamente cont\'{i}nuo
(sem vari\'{a}veis bin\'{a}rias ou inteiras) e cobre um horizonte de 15 dias.

\[
  \boxed{
    \begin{aligned}
      \min \; Z \;=\;\;
      & \underbrace{\sum_{p \in P}\sum_{l \in L_p}\sum_{i \in I}\sum_{t \in T}
          k_{i,p,l} \cdot x_{i,p,l,t}}_{\text{custo de produ\c{c}\~{a}o}}
        \;+\;
        \underbrace{\sum_{p \in P}\sum_{l \in L_p}\sum_{i \in I}\sum_{t \in T}
          v_{i,p,l} \cdot e_{i,p,l,t}}_{\text{custo de estoque}} \\[14pt]
      & +\;
        \underbrace{\sum_{i \in I}\sum_{o \in P}\sum_{\substack{d \in P \\ d \neq o}}\sum_{t \in T}
          ct_{o,d} \cdot \dfrac{f_{i,o,d,t}}{200{.}000}}_{\text{custo log\'{i}stico}}
        \;+\;
        \underbrace{\sum_{p \in P}\sum_{l \in L_p}\sum_{i \in I}\sum_{t \in T}
          M \cdot w_{i,p,l,t}}_{\text{penalidade por venda perdida}}
    \end{aligned}
  }
\]

\begin{center}
  \small\color{darkgray}
  Sujeito \`{a}s restri\c{c}\~{o}es R1--R9 detalhadas na Se\c{c}\~{a}o 7.
\end{center}

% ============================================================
\section{Conjuntos (\'Indices)}

\begin{center}
\renewcommand{\arraystretch}{1.4}
\begin{tabular}{>{\bfseries\color{ballblue}}c p{8cm} c r}
  \toprule
  S\'{i}mbolo & Descri\c{c}\~{a}o & \'{I}ndice & Cardinalidade \\
  \midrule
  $I$           & SKUs (produtos finais) & $i$ & 30 \\
  $P$           & Plantas (f\'{a}bricas) & $p,\,o,\,d$ & 10 \\
  $L_p$         & Linhas da planta $p$ & $l$ & 3 por planta \\
  $L$           & Todas as linhas ($L = \bigcup_{p \in P} L_p$) & $l$ & 30 \\
  $T$           & Per\'{i}odos $(t_1, \ldots, t_{15})$ & $t$ & 15 dias \\
  $F$           & Formatos de lata do portf\'{o}lio & $-$ & 3 \\
  $\mathcal{C}$ & Triplas cap\'{a}veis $(i,p,l)$: formato de $i$ = formato de $l \in L_p$ & $(i,p,l)$ & $\subseteq I \times L$ \\
  \bottomrule
\end{tabular}
\end{center}

\vspace{6pt}
\begin{mdframed}[style=destaque]
  \textbf{\color{ballblue}Conjunto $\mathcal{C}$ --- Capabilidade.}
  Os tr\^{e}s formatos s\~{a}o: \textit{9.1 Oz Sleek}, \textit{12 Oz Standard} e \textit{16 Oz Tall}.
  Cada \textbf{linha} $l$ opera com exatamente 1 formato, definido no arquivo de taxas
  (\texttt{sku\_size\_shape} em \texttt{input\_production\_rate}). Uma tripla $(i,p,l)$
  pertence a $\mathcal{C}$ se o formato do SKU $i$ coincide com o formato da linha $l$.
  Fora de $\mathcal{C}$, a produ\c{c}\~{a}o \'{e} proibida ($x_{i,p,l,t} = 0$ via R8).
  Os $\approx 10\%$ de pedidos incompat\'{i}veis s\~{a}o resolvidos via transfer\^{e}ncia $f_{i,o,d,t}$
  de uma planta apta, seguida de redistribui\c{c}\~{a}o interna $g_{i,p,l,t}$.
\end{mdframed}

% ============================================================
\section{Par\^{a}metros}

\subsection*{Log\'{i}sticos \textnormal{\small(\texttt{input\_logistic\_costs\_artificial.csv})}}

\begin{align*}
  cb_{o,d},\; cp_{o,d},\; cl_{o,d} &: \text{custos de combust\'{i}vel, ped\'{a}gio e m\~{a}o de obra da rota } o \to d \\
  ct_{o,d} &= cb_{o,d} + cp_{o,d} + cl_{o,d} \quad \text{(custo total por caminh\~{a}o)} \quad ct_{o,o} = 0
\end{align*}

\subsection*{Produ\c{c}\~{a}o \textnormal{\small(\texttt{input\_production\_rate} e \texttt{input\_sku\_costs})}}

\begin{align*}
  r_{p,l,t}   &: \text{taxa nominal de produ\c{c}\~{a}o da linha } l \in L_p \text{ no dia } t \text{ (latas/hora)} \\
               &\quad \text{\small Indexado por linha e dia --- n\~{a}o por SKU. Aus\^{e}ncia de registro implica produ\c{c}\~{a}o proibida.} \\[4pt]
  \rho        &: \text{efici\^{e}ncia operacional das m\'{a}quinas (OEE), } 0 < \rho \leq 1.
                  \text{ A taxa \emph{efetiva} \'{e} } r_{p,l,t}\cdot\rho. \text{ Base: } \rho = 1{,}0. \\[4pt]
  H           &: \text{janela di\'{a}ria de horas-m\'{a}quina por linha. Base: } H = 24\text{ h/dia}. \\[4pt]
  k_{i,p,l}   &: \text{custo de produ\c{c}\~{a}o do SKU } i \text{ na linha } l \in L_p \text{ (R\$/lata)} \\
  v_{i,p,l}   &: \text{custo de estocagem do SKU } i \text{ na linha } l \in L_p \text{ (R\$/lata/dia)}
\end{align*}

\subsection*{Demanda e armazenagem \textnormal{\small(\texttt{input\_production\_need} e \texttt{input\_line\_storage\_limit})}}

\begin{align*}
  \eta_{i,p,l,t}  &: \text{demanda do SKU } i \text{ na linha } l \in L_p \text{ com prazo no dia } t \text{ (latas)} \\
                   &\quad \text{\small Pode ser incompat\'{i}vel com o formato da planta (ru\'{i}do de $\approx$10\%).} \\[4pt]
  e^{\max}_{p,l}  &: \text{capacidade m\'{a}xima de armazenagem da linha } l \in L_p \text{ (latas) --- fixo.} \\
  e^{0}_{i,p,l}   &: \text{estoque inicial do SKU } i \text{ na linha } l \in L_p. \quad
                      \text{Neste estudo: } e^{0}_{i,p,l} = 0 \;\forall\, i,p,l.
\end{align*}

% ============================================================
\section{Fun\c{c}\~{a}o Auxiliar}

\[
  \phi : L \to P, \quad \phi(l) = p \qquad
  \text{(linha } l \mapsto \text{planta } p\text{)}
\]

Como $|L_p| = 3$ para todo $p$, a fun\c{c}\~{a}o $\phi$ \'{e} sobrejetora.
Sempre que $p$ e $l$ aparecem juntos nos \'{i}ndices, vale $p = \phi(l)$.

% ============================================================
\section{Vari\'{a}veis de Decis\~{a}o}

S\~{a}o sete vari\'{a}veis de decis\~{a}o, todas cont\'{i}nuas e n\~{a}o-negativas:

\vspace{6pt}
\begin{center}
\renewcommand{\arraystretch}{1.6}
\begin{tabular}{>{\bfseries\color{ballblue}}c p{9.5cm} l}
  \toprule
  Vari\'{a}vel & Descri\c{c}\~{a}o & Dom\'{i}nio \\
  \midrule
  $x_{i,p,l,t}$ & Volume \textbf{produzido} do SKU $i$ na linha $l \in L_p$, no dia $t$. S\'{o} definida para $(i,p,l) \in \mathcal{C}$. & $\mathbb{R}_+$ \\[4pt]
  $e_{i,p,l,t}$ & N\'{i}vel de \textbf{estoque} do SKU $i$ na linha $l \in L_p$ ao final do dia $t$. & $\mathbb{R}_+$ \\[4pt]
  $a_{i,p,l,t}$ & Volume de \textbf{atendimento local} do SKU $i$ a partir da linha $l \in L_p$ no dia $t$. & $\mathbb{R}_+$ \\[4pt]
  $f_{i,o,d,t}$ & Volume \textbf{transferido} do SKU $i$ da planta $o$ para a planta $d \neq o$, no dia $t$. & $\mathbb{R}_+$ \\[4pt]
  $g_{i,p,l,t}$ & Volume \textbf{redistribu\'{i}do internamente} da doca de entrada da planta $p$ para o armazém da linha $l$, no dia $t$. & $\mathbb{R}_+$ \\[4pt]
  $s_{i,p,l,t}$ & Volume \textbf{retirado do armazém} da linha $l \in L_p$ para a doca de sa\'{i}da da planta $p$, no dia $t$. & $\mathbb{R}_+$ \\[4pt]
  $w_{i,p,l,t}$ & Volume de \textbf{demanda n\~{a}o atendida} (venda perdida) do SKU $i$ na linha $l \in L_p$, no dia $t$. Vari\'{a}vel artificial de folga. & $\mathbb{R}_+$ \\
  \bottomrule
\end{tabular}
\end{center}

\vspace{10pt}

\begin{mdframed}[style=destaque]
  \textbf{\color{ballblue}Por que $f_{i,o,d,t}$ e n\~{a}o $x$ com destino?}

  \medskip
  O frete \'{e} pago pelo \textbf{movimento f\'{i}sico entre plantas}, n\~{a}o pela produ\c{c}\~{a}o.
  Se Jacare\'{i} produz 1.000.000 de latas e guarda tudo localmente por 3 dias,
  nenhum caminh\~{a}o \'{e} usado --- custo log\'{i}stico zero.
  O caminho f\'{i}sico das latas \'{e}:
  \[
    x_{i,p,l,t}
    \xrightarrow{\text{estoque}}
    e_{i,p,l,t}
    \xrightarrow{s: \text{ doca de sa\'{i}da}}
    f_{i,p,d,t}
    \xrightarrow{g: \text{ doca de entrada}}
    e_{i,d,l',t}
    \xrightarrow{a}
    \text{cliente}
  \]
  Cada vari\'{a}vel carrega seu pr\'{o}prio custo:
  $x \to k_{i,p,l}$,\quad $e \to v_{i,p,l}$,\quad $f \to ct_{o,d}$.
\end{mdframed}

\vspace{8pt}

\begin{mdframed}[style=destaque]
  \textbf{\color{ballblue}O n\'{o} de transbordo: por que $g$ e $s$?}

  \medskip
  Caminh\~{o}es chegam e saem pelo p\'{a}tio da f\'{a}brica (a doca da planta $p$),
  n\~{a}o diretamente na porta de cada linha. As vari\'{a}veis $g$ e $s$ modelam
  as \textbf{empilhadeiras} que fazem a distribui\c{c}\~{a}o interna:

  \begin{itemize}[leftmargin=1.5em, itemsep=2pt]
    \item $g_{i,p,l,t}$: caminh\~{a}o descarrega no p\'{a}tio $\to$ empilhadeira leva ao armazém da linha $l$.
    \item $s_{i,p,l,t}$: empilhadeira retira do armazém da linha $l$ $\to$ leva \`{a} doca para encher caminh\~{a}o.
  \end{itemize}

  \medskip
  \textbf{Por que n\~{a}o usar $f$ direto no balan\c{c}o de estoque?}
  Se sum\'{a}ssemos $\sum_o f_{i,o,p,t}$ direto na equa\c{c}\~{a}o de cada linha $l$,
  o solver \textit{clonaria} as latas recebidas para todos os armazéns simultaneamente.
  As restri\c{c}\~{o}es R6 e R7 garantem o \textbf{Princ\'{i}pio da Conserva\c{c}\~{a}o de Massa}
  no transbordo: o que entra no p\'{a}tio sai exatamente para os armazéns, e vice-versa.
\end{mdframed}

\vspace{8pt}

\begin{mdframed}[style=destaque]
  \textbf{\color{ballblue}Como o estoque $e_{i,p,l,t}$ se forma --- balan\c{c}o intuitivo}

  \medskip
  \[
    \begin{aligned}
      e_{i,p,l,t} \;=\;\;
      & \underbrace{e_{i,p,l,\,t-1}}_{\text{estoque anterior}}
        \;+\; \underbrace{x_{i,p,l,t}}_{\text{produ\c{c}\~{a}o}}
        \;+\; \underbrace{g_{i,p,l,t}}_{\text{entrada da doca}} \\[10pt]
      & -\; \underbrace{s_{i,p,l,t}}_{\text{sa\'{i}da para doca}}
        \;-\; \underbrace{a_{i,p,l,t}}_{\text{atendimento local}}
    \end{aligned}
  \]
  \textit{Formalizado como restri\c{c}\~{a}o de igualdade em R2, com condi\c{c}\~{a}o inicial $e_{i,p,l,0} = e^0_{i,p,l} = 0$.}
\end{mdframed}

\vspace{8pt}

\begin{mdframed}[style=destaque]
  \textbf{\color{ballblue}O que \'{e} $w_{i,p,l,t}$ e por que ela existe? --- A l\'{o}gica do Big-$M$}

  \medskip
  No mundo real, quando a demanda excede a capacidade f\'{i}sica, ocorre \textbf{venda perdida}
  --- n\~{a}o o colapso do universo. Para evitar que o modelo declare \textit{infeasible}
  em cen\'{a}rios de estresse (ex: $150\%$ de ocupa\c{c}\~{a}o), introduzimos $w$ como
  uma \textbf{vari\'{a}vel artificial de folga}.

  \medskip
  \textbf{O balan\c{c}o econ\^{o}mico (Big-$M$):} Se $w$ custasse zero, o solver
  \textit{deixaria de produzir} e jogaria toda a demanda em $w$ (pois produzir custa dinheiro).
  Punimos $w$ com um custo $M = R\$\,1.000.000$ por lata n\~{a}o entregue --- t\~{a}o alto que
  o modelo s\'{o} recorre a $w$ quando esgota todas as possibilidades f\'{i}sicas.
  Esse mecanismo \'{e} conhecido em Pesquisa Operacional como \textbf{penalidade Big-$M$}.
\end{mdframed}

% ============================================================
\section{Guia das Vari\'{a}veis: o que cada uma representa fisicamente}

Esta se\c{c}\~{a}o \'{e} destinada a quem conhece Programa\c{c}\~{a}o Linear mas n\~{a}o
conhece a opera\c{c}\~{a}o de uma malha de f\'{a}bricas de latas. O objetivo \'{e} tornar
claro \textbf{por que} cada vari\'{a}vel existe e o que ela representa no ch\~{a}o de f\'{a}brica
--- antes de v\^{e}-la numa equa\c{c}\~{a}o.

% ── a ────────────────────────────────────────────────────────────────────────
\vspace{6pt}
\begin{mdframed}[style=destaque]
\textbf{\color{ballblue}$a_{i,p,l,t}$ --- Atendimento local: a \'{u}nica vari\'{a}vel que toca o cliente}

\medskip
No modelo, \textbf{nenhuma outra vari\'{a}vel entrega lata na m\~{a}o do cliente}.
$x$ produz, $e$ guarda, $f$ transporta, $g$ distribui internamente, $s$ coleta para expedir ---
mas nenhuma delas encerra o pedido. Quem encerra \'{e} $a_{i,p,l,t}$.

\medskip
Pense assim: um pedido chega na linha JAC\_L2 exigindo 50.000 latas de Brahma\_Lata at\'{e}
o dia 5. O pedido s\'{o} \'{e} considerado atendido quando $a_{\text{Brahma},\text{JAC},\text{JAC\_L2},5}
\geq 50.000$. O caminho que essas latas percorreram at\'{e} chegar l\'{a} --- produzidas
localmente, chegadas de Extrema via transfer\^{e}ncia, ou acumuladas desde o dia 3 ---
n\~{a}o importa para o cliente. Importa que $a$ seja suficiente.

\medskip
\textbf{Por que $a$ \'{e} indexado por linha e n\~{a}o por planta?}
Porque a demanda $\eta_{i,p,l,t}$ chegou amarrada a uma linha espec\'{i}fica
(campo \texttt{production\_line} no CSV). O atendimento precisa ser rastreado
na mesma granularidade do pedido para que a restri\c{c}\~{a}o R3 ($a + w = \eta$)
fa\c{c}a sentido f\'{i}sico.

\medskip
\textbf{O que $a$ n\~{a}o \'{e}:} $a$ n\~{a}o \'{e} produ\c{c}\~{a}o. Uma linha pode atender
um pedido inteiramente com latas que vieram de outra planta via $f \to g$,
sem produzir nada localmente ($x = 0$). Nesse caso, $a > 0$ mas $x = 0$ --- algo
imposs\'{i}vel de modelar se confundirmos as duas vari\'{a}veis.
\end{mdframed}

% ── e ────────────────────────────────────────────────────────────────────────
\vspace{8pt}
\begin{mdframed}[style=destaque]
\textbf{\color{ballblue}$e_{i,p,l,t}$ --- Por que o estoque \'{e} da \textbf{linha}, n\~{a}o da planta?}

\medskip
Quem conhece modelos de rede esperaria que o estoque fosse indexado pela planta $p$
--- afinal, o armazém \'{e} f\'{i}sico e pertence \`{a} f\'{a}brica. Mas aqui o estoque
\'{e} indexado pela \textbf{linha} $l$, e h\'{a} uma raz\~{a}o precisa para isso.

\medskip
A demanda $\eta_{i,p,l,t}$ chega amarrada a uma linha espec\'{i}fica. O atendimento
$a_{i,p,l,t}$ tamb\'{e}m \'{e} por linha. Para que o balan\c{c}o de estoque R2 feche
corretamente ($e_{t} = e_{t-1} + x + g - s - a$), o estoque precisa estar
na mesma granularidade de $a$ e $x$ --- ou seja, por linha.

\medskip
\textbf{Analogia:} imagine que cada linha \'{e} um corredor separado de supermercado.
O cliente que pediu latas no corredor A n\~{a}o aceita receber do corredor B sem
antes algum funcion\'{a}rio transferir a mercadoria entre corredores
(isso \'{e} o $g$ --- a redistribui\c{c}\~{a}o interna).
\end{mdframed}

% ── g e s ────────────────────────────────────────────────────────────────────
\vspace{8pt}
\begin{mdframed}[style=destaque]
\textbf{\color{ballblue}$g_{i,p,l,t}$ e $s_{i,p,l,t}$ --- As empilhadeiras que ningu\'{e}m vê}

\medskip
Essas s\~{a}o as vari\'{a}veis mais ``escondidas'' do modelo. Elas existem por uma
necessidade matem\'{a}tica clara: $f$ \'{e} indexado por \textbf{planta} ($o \to d$),
mas o estoque \'{e} indexado por \textbf{linha} ($l$ dentro de $p$). Precisa haver
algo que faça a ponte entre os dois n\'{i}veis hier\'{a}rquicos.

\medskip
\begin{center}
\renewcommand{\arraystretch}{1.3}
\small
\begin{tabular}{>{\bfseries\color{ballblue}}c p{4cm} p{5.5cm}}
  \toprule
  Vari\'{a}vel & O que representa & O que aconteceria sem ela \\
  \midrule
  $g_{i,p,l,t}$ &
    Latas descarregadas no p\'{a}tio de $p$ levadas ao armazém da linha $l$ &
    O solver ``clonaria'' o frete recebido para \textbf{todas} as linhas ao mesmo
    tempo, multiplicando latas do nada \\[6pt]
  $s_{i,p,l,t}$ &
    Latas retiradas do armazém da linha $l$ levadas \`{a} doca de expedi\c{c}\~{a}o de $p$ &
    N\~{a}o haveria como saber \textbf{de qual linha} sa\'{i}ram as latas expedidas,
    impossibilitando o balan\c{c}o por linha \\
  \bottomrule
\end{tabular}
\end{center}

\medskip
As restri\c{c}\~{o}es R6 e R7 s\~{a}o as ``cancelas'' que garantem conserva\c{c}\~{a}o de massa:
tudo que chega via $f$ deve sair via $g$; tudo que sai via $f$ deve ter sido coletado via $s$.
Sem essas cancelas, o modelo poderia \textit{criar} ou \textit{destruir} latas no transbordo.
\end{mdframed}

% ── f ────────────────────────────────────────────────────────────────────────
\vspace{8pt}
\begin{mdframed}[style=destaque]
\textbf{\color{ballblue}$f_{i,o,d,t}$ --- Transfer\^{e}ncia: quando e por que o modelo decide usá-la?}

\medskip
$f$ existe para resolver dois problemas do mundo real:

\begin{enumerate}[leftmargin=1.5em, itemsep=4pt]
  \item \textbf{O ru\'{i}do de formato ($\approx 10\%$ dos pedidos):}
    quando um pedido chega a uma linha que n\~{a}o sabe produzir aquele formato
    (ex: pedido de \textit{9.1 Oz Sleek} para uma linha de \textit{12 Oz Standard}),
    a \'{u}nica sa\'{i}da do modelo \'{e} buscar outra planta apta e transferir.
    O solver descobre isso sozinho: $x = 0$ por R8, ent\~{a}o para atender $a \geq \eta$
    via R3, precisa de $g > 0$, que exige $f > 0$ via R6.

  \item \textbf{Arbitragem de capacidade:}
    mesmo sem ru\'{i}do, pode ser \'{o}timo produzir num lugar mais barato e transferir.
    Se a planta de Extrema tem custo de produ\c{c}\~{a}o menor e folga produtiva,
    o modelo pode decidir produzir l\'{a} e mandar caminh\~{o}es para Jacare\'{i},
    desde que o custo do frete ($ct_{o,d}$) n\~{a}o anule a economia.
\end{enumerate}

\medskip
\textbf{Por que $f$ n\~{a}o \'{e} indexado por linha?}
O caminh\~{a}o n\~{a}o sabe (nem precisa saber) para qual linha dentro da planta destino
as latas v\~{a}o. Ele descarrega no p\'{a}tio e as empilhadeiras ($g$) decidem.
Indexar $f$ por linha forçaria a origem a conhecer a organiza\c{c}\~{a}o interna do destino
--- uma informa\c{c}\~{a}o desnecessária e acoplamento errado entre plantas.
\end{mdframed}

% ── w ────────────────────────────────────────────────────────────────────────
\vspace{8pt}
\begin{mdframed}[style=destaque]
\textbf{\color{ballblue}$w_{i,p,l,t}$ --- Venda perdida: a vari\'{a}vel que evita o \textit{infeasible}}

\medskip
Em cen\'{a}rios onde a demanda excede a capacidade f\'{i}sica total da malha,
um modelo sem $w$ simplesmente \textbf{n\~{a}o tem solu\c{c}\~{a}o} --- ele declara
\textit{infeasible} e para. Isso \'{e} um problema porque:

\begin{itemize}[leftmargin=1.5em, itemsep=2pt]
  \item A inviabilidade n\~{a}o diz \textit{o quanto} a malha ficou devendo.
  \item N\~{a}o \'{e} poss\'{i}vel comparar cen\'{a}rios: um cen\'{a}rio invi\'{a}vel
        n\~{a}o tem custo \'{o}timo para comparar com outro.
  \item N\~{a}o \'{e} poss\'{i}vel fazer an\'{a}lise de sensibilidade.
\end{itemize}

\medskip
$w$ transforma o problema invi\'{a}vel em \textbf{sempre vi\'{a}vel}: qualquer demanda
pode ser ``atendida'' por $a$ ou ``absorvida'' por $w$. O truque \'{e} a penalidade
Big-$M = R\$\,10^6$ por lata em $w$ --- t\~{a}o cara que o solver s\'{o} recorre a
ela quando esgota todo o potencial produtivo e log\'{i}stico da rede.

\medskip
\textbf{Como ler os resultados:} se \texttt{custo\_penalidade} = 0 no output,
a malha atendeu 100\% da demanda. Se for positivo, $w > 0$ em alguma linha/dia
--- a raz\~{a}o \'{e} f\'{i}sica (capacidade insuficiente), n\~{a}o matem\'{a}tica.
\end{mdframed}

% ============================================================
\section{Fun\c{c}\~{a}o Objetivo}

\[
  \begin{aligned}
    \min \; Z \;=\;\;
    & \sum_{p \in P}\sum_{\substack{l \in L_p \\ (i,p,l)\,\in\,\mathcal{C}}}\sum_{i \in I}\sum_{t \in T}
        k_{i,p,l} \cdot x_{i,p,l,t}
      \;+\;
      \sum_{p \in P}\sum_{l \in L_p}\sum_{i \in I}\sum_{t \in T}
        v_{i,p,l} \cdot e_{i,p,l,t} \\[10pt]
    & +\;
      \sum_{i \in I}\sum_{o \in P}\sum_{\substack{d \in P \\ d \,\neq\, o}}\sum_{t \in T}
        ct_{o,d} \cdot \frac{f_{i,o,d,t}}{200{.}000}
      \;+\;
      \sum_{p \in P}\sum_{l \in L_p}\sum_{i \in I}\sum_{t \in T}
        M \cdot w_{i,p,l,t}
  \end{aligned}
\]

\begin{center}
\renewcommand{\arraystretch}{1.3}
\small
\begin{tabular}{>{\bfseries\color{ballblue}}c l l}
  \toprule
  Termo & F\'{o}rmula & Interpreta\c{c}\~{a}o \\
  \midrule
  Produ\c{c}\~{a}o  & $k_{i,p,l} \cdot x_{i,p,l,t}$                     & Custo unit\'{a}rio (R\$/lata) $\times$ volume produzido \\
  Estoque        & $v_{i,p,l} \cdot e_{i,p,l,t}$                     & Custo de carregar 1 lata no armazém por 1 dia \\
  Log\'{i}stica  & $ct_{o,d} \cdot \frac{f_{i,o,d,t}}{200{.}000}$   & Custo por caminh\~{a}o $\times$ n\'{u}mero de caminh\~{o}es \\
  Venda perdida  & $M \cdot w_{i,p,l,t}$                              & Penalidade Big-$M$ por demanda n\~{a}o atendida \\
  \bottomrule
\end{tabular}
\end{center}

\begin{center}
  \small\color{darkgray}
  Capacidade do caminh\~{a}o: $25\text{ pallets} \times 8.000\text{ latas} = 200.000\text{ latas/caminh\~{a}o}$
\end{center}

% ============================================================
\newpage
\section{Restri\c{c}\~{o}es}

\begin{mdframed}[style=destaque]
\textbf{\color{ballblue}R1 --- Limite de armazenagem.}
\medskip

O espa\c{c}o f\'{i}sico do armazém \'{e} compartilhado entre todos os SKUs.
A soma do estoque de todos os produtos n\~{a}o pode ultrapassar a capacidade da linha:
\[
  \sum_{i \in I} e_{i,p,l,t} \;\leq\; e^{\max}_{p,l}
  \qquad \forall\, p \in P,\; l \in L_p,\; t \in T
\]
\end{mdframed}

\vspace{8pt}

\begin{mdframed}[style=destaque]
\textbf{\color{ballblue}R2 --- Balan\c{c}o de estoque.}
\medskip

Equa\c{c}\~{a}o de conserva\c{c}\~{a}o de massa do armazém da linha. Unificada usando
$e_{i,p,l,0} = e^0_{i,p,l} = 0$:
\[
  e_{i,p,l,t} = e_{i,p,l,t-1} + x_{i,p,l,t} + g_{i,p,l,t} - s_{i,p,l,t} - a_{i,p,l,t}
  \qquad \forall\, i \in I,\; p \in P,\; l \in L_p,\; t \in T
\]
Esta \'{e} uma igualdade estrita que \textbf{acopla os 15 per\'{i}odos}: o estoque de hoje
\'{e} a condi\c{c}\~{a}o inicial de amanh\~{a}. O frete ($f$) n\~{a}o aparece aqui
diretamente --- ele entra via $g$ (entrada) e sai via $s$ (sa\'{i}da), garantindo
que a redistribui\c{c}\~{a}o interna \'{e} tratada pelas docas (R6 e R7).
\end{mdframed}

\vspace{8pt}

\begin{mdframed}[style=destaque]
\textbf{\color{ballblue}R3 --- Atendimento \`{a} demanda.}
\medskip

O atendimento local $a$ somado \`{a} venda perdida $w$ deve igualar exatamente a demanda:
\[
  a_{i,p,l,t} + w_{i,p,l,t} = \eta_{i,p,l,t}
  \qquad \forall\, i \in I,\; p \in P,\; l \in L_p,\; t \in T
\]
O frete e a produ\c{c}\~{a}o \textbf{n\~{a}o aparecem aqui} para evitar dupla contagem:
eles j\'{a} alimentaram o estoque via R2, e o estoque alimenta $a$.
O solver s\'{o} aciona $w$ quando esgota os 24\,h produtivos de todas as linhas.
\end{mdframed}

\vspace{8pt}

\begin{mdframed}[style=destaque]
\textbf{\color{ballblue}R4 --- Capacidade produtiva (janela de horas-m\'{a}quina).}
\medskip

A soma das horas consumidas por todos os SKUs na linha $l$ n\~{a}o pode ultrapassar a
janela di\'{a}ria $H$. O tempo para produzir $x_{i,p,l,t}$ latas depende da taxa nominal
$r_{p,l,t}$ ajustada pela efici\^{e}ncia operacional $\rho$ (OEE): a taxa efetiva \'{e}
$r_{p,l,t}\cdot\rho$, logo o tempo \'{e} $x_{i,p,l,t} / (r_{p,l,t}\cdot\rho)$ horas:
\[
  \sum_{i \in I} \frac{x_{i,p,l,t}}{r_{p,l,t}\cdot\rho} \;\leq\; H_{p}
  \qquad \forall\, p \in P,\; l \in L_p,\; t \in T
\]
onde $H_{p} = H = 24$\,h para plantas operacionais e $H_{p} = 0$ quando a planta $p$
est\'{a} desativada (par\^{a}metro \texttt{fabrica\_quebrada} da an\'{a}lise de
sensibilidade). Na configura\c{c}\~{a}o base, $\rho = 1{,}0$ e $H_{p} = 24$\,h.
Para $(i,p,l) \notin \mathcal{C}$, temos $x_{i,p,l,t} = 0$ por R8, logo esses
termos n\~{a}o contribuem para a soma.
\end{mdframed}

\vspace{8pt}

\begin{mdframed}[style=destaque]
\textbf{\color{ballblue}R5 --- Capacidade log\'{i}stica de expedi\c{c}\~{a}o.}
\medskip

Cada planta pode expedir no m\'{a}ximo 80 caminh\~{o}es por dia (sa\'{i}da),
somando todos os SKUs e destinos:
\[
  \sum_{i \in I} \sum_{\substack{d \in P \\ d \neq p}}
  \frac{f_{i,p,d,t}}{200{.}000} \;\leq\; 80
  \qquad \forall\, p \in P,\; t \in T
\]
\end{mdframed}

\vspace{8pt}

\begin{mdframed}[style=destaque]
\textbf{\color{ballblue}R6 e R7 --- Conserva\c{c}\~{a}o de massa no transbordo (docas).}
\medskip

Tudo que os caminh\~{o}es descarregam no p\'{a}tio deve ser exatamente distribu\'{i}do
pelos armazéns das linhas (R6), e tudo que os armazéns enviam ao p\'{a}tio deve
ser exatamente carregado nos caminh\~{o}es (R7):

\[
  \sum_{l \in L_p} g_{i,p,l,t} = \sum_{\substack{o \in P \\ o \neq p}} f_{i,o,p,t}
  \qquad \forall\, i \in I,\; p \in P,\; t \in T \tag{R6}
\]
\[
  \sum_{l \in L_p} s_{i,p,l,t} = \sum_{\substack{d \in P \\ d \neq p}} f_{i,p,d,t}
  \qquad \forall\, i \in I,\; p \in P,\; t \in T \tag{R7}
\]
\end{mdframed}

\vspace{8pt}

\begin{mdframed}[style=destaque]
\textbf{\color{ballblue}R8 --- Capabilidade de formato.}
\medskip

Linhas s\'{o} podem produzir SKUs compat\'{i}veis com seu formato:
\[
  x_{i,p,l,t} = 0 \qquad \forall\, (i,p,l) \notin \mathcal{C},\; t \in T
\]
\end{mdframed}

\vspace{8pt}

\begin{mdframed}[style=destaque]
\textbf{\color{ballblue}R9 --- N\~{a}o-negatividade.}
\[
  x_{i,p,l,t},\; e_{i,p,l,t},\; a_{i,p,l,t},\; f_{i,o,d,t},\;
  g_{i,p,l,t},\; s_{i,p,l,t},\; w_{i,p,l,t} \;\geq\; 0
\]
\end{mdframed}

% ============================================================
\newpage
\section{Formula\c{c}\~{a}o Compacta do PL}

\begin{mdframed}[backgroundcolor=blue!5, linecolor=ballblue, linewidth=2pt,
  innertopmargin=14pt, innerbottommargin=14pt,
  innerleftmargin=14pt, innerrightmargin=14pt]

\textbf{\color{ballblue}min}
\[
  \begin{aligned}
    Z \;=\;\;
    & \sum_{p,l,i,t} k_{i,p,l} \cdot x_{i,p,l,t}
      \;+\; \sum_{p,l,i,t} v_{i,p,l} \cdot e_{i,p,l,t}
      \;+\; \sum_{i,o,d,t} ct_{o,d} \cdot \tfrac{f_{i,o,d,t}}{200{.}000}
      \;+\; \sum_{p,l,i,t} M \cdot w_{i,p,l,t}
  \end{aligned}
\]

\textbf{\color{ballblue}s.a.}
\begin{alignat}{3}
  &\textstyle\sum_{i} e_{i,p,l,t}
    &&\leq e^{\max}_{p,l}
    &&\quad\forall\,p,l,t \tag{R1}\\[6pt]
  &e_{i,p,l,t} = e_{i,p,l,t-1} + x_{i,p,l,t} + g_{i,p,l,t} - s_{i,p,l,t} - a_{i,p,l,t}
    &&
    &&\quad\forall\,i,p,l,t \tag{R2}\\[6pt]
  &a_{i,p,l,t} + w_{i,p,l,t}
    &&= \eta_{i,p,l,t}
    &&\quad\forall\,i,p,l,t \tag{R3}\\[6pt]
  &\textstyle\sum_{i} \tfrac{x_{i,p,l,t}}{r_{p,l,t}\cdot\rho}
    &&\leq H_{p}
    &&\quad\forall\,p,l,t \tag{R4}\\[6pt]
  &\textstyle\sum_{i,d\neq p} \tfrac{f_{i,p,d,t}}{200{.}000}
    &&\leq 80
    &&\quad\forall\,p,t \tag{R5}\\[6pt]
  &\textstyle\sum_{l \in L_p} g_{i,p,l,t}
    &&= \textstyle\sum_{o \neq p} f_{i,o,p,t}
    &&\quad\forall\,i,p,t \tag{R6}\\[6pt]
  &\textstyle\sum_{l \in L_p} s_{i,p,l,t}
    &&= \textstyle\sum_{d \neq p} f_{i,p,d,t}
    &&\quad\forall\,i,p,t \tag{R7}\\[6pt]
  &x_{i,p,l,t} = 0
    &&
    &&\quad\forall\,(i,p,l)\notin\mathcal{C},\,t \tag{R8}\\[6pt]
  &x,e,a,f,g,s,w \geq 0
    &&
    &&\quad\forall\,i,p,l,o,d,t \tag{R9}
\end{alignat}
\end{mdframed}

\vspace{8pt}
{\color{accentgray}\rule{\linewidth}{0.5pt}}
\begin{center}
  \small\color{darkgray}
  $6 \times 30 \times 10 \times 3 \times 15 = \mathbf{81{.}000}$ vari\'{a}veis operacionais
  $+\; 30 \times 100 \times 15 = \mathbf{45{.}000}$ vari\'{a}veis de transfer\^{e}ncia
  $= \approx \mathbf{126{.}000}$ vari\'{a}veis no total.
\end{center}

% ============================================================
\newpage
\section{Implementa\c{c}\~{a}o em Julia/JuMP}

Esta se\c{c}\~{a}o explica como o modelo matem\'{a}tico se traduz em c\'{o}digo, passo a passo.
O objetivo \'{e} que mesmo quem nunca programou em Julia consiga entender a l\'{o}gica por tr\'{a}s de cada bloco.

% ── 9.1 ─────────────────────────────────────────────────────────────────────
\subsection*{9.1 Pacotes utilizados}

\begin{mdframed}[style=codigo]
\begin{lstlisting}
using JuMP      # linguagem de modelagem matematica (escreve o LP)
using Gurobi    # solver comercial que resolve o LP
using CSV       # leitura de arquivos .csv
using DataFrames  # estrutura de tabela (como um Excel no codigo)
using Printf    # formatacao de numeros no terminal
using Dates     # data/hora para nomear os arquivos de saida
\end{lstlisting}
\end{mdframed}

\begin{mdframed}[style=destaque]
\textbf{\color{ballblue}O que \'{e} JuMP?} JuMP \'{e} uma linguagem de modelagem embutida no Julia.
Em vez de implementar o simplex manualmente, voc\^{e} \textbf{descreve} o modelo
(vari\'{a}veis, restri\c{c}\~{o}es, objetivo) e o JuMP monta a matriz LP e entrega ao solver.
O Gurobi \'{e} o solver --- o motor matem\'{a}tico que de fato encontra a solu\c{c}\~{a}o \'{o}tima.
\end{mdframed}

% ── 9.2 ─────────────────────────────────────────────────────────────────────
\subsection*{9.2 Leitura e prepara\c{c}\~{a}o dos dados}

\begin{mdframed}[style=codigo]
\begin{lstlisting}
# Extrai o dia da string de data (ex: "2024-01-05" -> 5)
function extrair_dia(data_str::String)
    return parse(Int, string(data_str)[9:10])
end

# Converte numero brasileiro "1.234,56" para Float64 1234.56
function parse_br_float(val)
    ismissing(val) || val == "" && return 0.0
    return parse(Float64, replace(String(val), "," => "."))
end
\end{lstlisting}
\end{mdframed}

\begin{mdframed}[style=destaque]
\textbf{\color{ballblue}Por que essas fun\c{c}\~{o}es?}
Os CSVs usam o padr\~{a}o brasileiro de n\'{u}meros (v\'{i}rgula como decimal).
O Julia espera ponto. Sem a convers\~{a}o, \texttt{parse} lan\c{c}aria um erro.
A extra\c{c}\~{a}o do dia resolve o fato de que o campo \texttt{deadline} vem
como data completa (\texttt{"2024-01-05"}), mas o modelo usa $t \in \{1,\ldots,15\}$.
\end{mdframed}

% ── 9.3 ─────────────────────────────────────────────────────────────────────
\subsection*{9.3 Extra\c{c}\~{a}o dos conjuntos}

\begin{mdframed}[style=codigo]
\begin{lstlisting}
I   = unique(df_need.product_id)   # conjunto I: 30 SKUs
P   = unique(df_rate.plant)        # conjunto P: 10 plantas
T   = 1:15                         # conjunto T: dias 1 a 15

# L_p: dicionario planta -> lista de linhas daquela planta
L_p = Dict(p => unique(df_rate[df_rate.plant .== p,
           :production_line]) for p in P)

# Conjunto C: triplas (i, p, l) onde o formato do SKU bate
# com o formato que a linha sabe produzir
formatos_da_linha = Dict(r.production_line => r.sku_size_shape
                         for r in eachrow(df_rate))
C = Set(
    (r.product_id, r.plant, r.production_line)
    for r in eachrow(df_need)
    if haskey(formatos_da_linha, r.production_line) &&
       r.sku_size_shape == formatos_da_linha[r.production_line]
)
\end{lstlisting}
\end{mdframed}

\begin{mdframed}[style=destaque]
\textbf{\color{ballblue}O \texttt{Dict} como $L_p$ e o \texttt{Set} como $\mathcal{C}$:}
O dicion\'{a}rio \texttt{L\_p} implementa $L_p$ diretamente --- iterar
\texttt{for p in P, l in L\_p[p]} \'{e} exatamente $\sum_{p \in P}\sum_{l \in L_p}$.
O \texttt{Set} $C$ armazena as triplas cap\'{a}veis e permite verifica\c{c}\~{a}o
\texttt{(i,p,l) \ensuremath{\in} C} em tempo $O(1)$ por hash, cr\'{i}tico dado que ocorre
em cada itera\c{c}\~{a}o dos somatórios.
\end{mdframed}

% ── 9.4 ─────────────────────────────────────────────────────────────────────
\subsection*{9.4 Constru\c{c}\~{a}o dos par\^{a}metros}

\begin{mdframed}[style=codigo]
\begin{lstlisting}
# parse_br_float aplicado em TODOS os campos numericos dos CSVs
k     = Dict((r.product_id, r.plant, r.production_line) =>
             parse_br_float(r.production_cost) for r in eachrow(df_costs))
v     = Dict((r.product_id, r.plant, r.production_line) =>
             parse_br_float(r.inventory_cost)  for r in eachrow(df_costs))
ct    = Dict((r.origem, r.destino) =>
             parse_br_float(r.custo_total_frete) for r in eachrow(df_log))
e_max = Dict((r.plant, r.production_line) =>
             parse_br_float(r.storage_limit) for r in eachrow(df_limit))

# Demanda: inicializa com 0 para cobrir combinacoes sem pedido no CSV
eta = Dict{Tuple{Any,Any,Any,Int}, Float64}()
for i in I, p in P, l in L_p[p], t in T
    eta[(i, p, l, t)] = 0.0
end
for r in eachrow(df_need)
    t = extrair_dia(string(r.deadline))
    chave = (r.product_id, r.plant, r.production_line, t)
    haskey(eta, chave) && (eta[chave] += r.prod_need)
end
\end{lstlisting}
\end{mdframed}

\begin{mdframed}[style=destaque]
\textbf{\color{ballblue}Por que \texttt{parse\_br\_float} em todos os par\^{a}metros?}
A vers\~{a}o anterior s\'{o} convertia a coluna \texttt{prod\_need}.
Os campos \texttt{production\_cost}, \texttt{inventory\_cost}, \texttt{custo\_total\_frete}
e \texttt{storage\_limit} tamb\'{e}m podem usar v\'{i}rgula decimal nos CSVs brasileiros.
Sem a convers\~{a}o, esses valores seriam lidos como \texttt{String} e causariam erros
silenciosos ou crash na montagem do modelo.
\end{mdframed}

% ── 9.5 ─────────────────────────────────────────────────────────────────────
\subsection*{9.5 Vari\'{a}veis de decis\~{a}o no JuMP}

\begin{mdframed}[style=codigo]
\begin{lstlisting}
modelo = Model(Gurobi.Optimizer)
set_silent(modelo)  # suprime o log do Gurobi

@variable(modelo, x[i in I, p in P, l in L_p[p], t in T] >= 0)
@variable(modelo, e[i in I, p in P, l in L_p[p], t in T] >= 0)
@variable(modelo, a[i in I, p in P, l in L_p[p], t in T] >= 0)
@variable(modelo, f[i in I, o in P, d in P, t in T; o != d] >= 0)
@variable(modelo, g[i in I, p in P, l in L_p[p], t in T] >= 0)
@variable(modelo, s[i in I, p in P, l in L_p[p], t in T] >= 0)
@variable(modelo, w[i in I, p in P, l in L_p[p], t in T] >= 0)
\end{lstlisting}
\end{mdframed}

\begin{mdframed}[style=destaque]
\textbf{\color{ballblue}Como ler \texttt{@variable}:}
Cada \texttt{@variable} cria automaticamente \textbf{uma vari\'{a}vel por combina\c{c}\~{a}o}
de \'{i}ndices. Por exemplo, \texttt{x[i in I, p in P, l in L\_p[p], t in T]}
cria $30 \times 10 \times 3 \times 15 = 13.500$ vari\'{a}veis.
O \texttt{o != d} em \texttt{f} garante que transferência para a mesma planta n\~{a}o existe ---
implementando diretamente a condi\c{c}\~{a}o $o \neq d$ do modelo.
O \texttt{>= 0} implementa R9 automaticamente.
\end{mdframed}

% ── 9.6 ─────────────────────────────────────────────────────────────────────
\subsection*{9.6 Fun\c{c}\~{a}o objetivo (com $\mathcal{C}$ expl\'{i}cito)}

\begin{mdframed}[style=codigo]
\begin{lstlisting}
@objective(modelo, Min,
    # Producao: so triplas em C contribuem
    sum(k[(i,p,l)] * x[i,p,l,t]
        for p in P, l in L_p[p], i in I, t in T
        if (i,p,l) in C) +
    # Estoque: todos os armazens
    sum(v[(i,p,l)] * e[i,p,l,t]
        for p in P, l in L_p[p], i in I, t in T
        if haskey(v,(i,p,l))) +
    # Logistica: custo por caminhao
    sum(ct[(o,d)] * (f[i,o,d,t] / 200_000)
        for i in I, o in P, d in P, t in T if o != d) +
    # Big-M: penalidade por venda perdida
    sum(1_000_000 * w[i,p,l,t]
        for p in P, l in L_p[p], i in I, t in T)
)
\end{lstlisting}
\end{mdframed}

\begin{mdframed}[style=destaque]
\textbf{\color{ballblue}Conjunto $\mathcal{C}$ vs \texttt{haskey}:}
A vers\~{a}o anterior usava \texttt{haskey(k,(i,p,l))} como proxy para $\mathcal{C}$.
O c\'{o}digo correto usa \texttt{(i,p,l) in C}, onde $C$ foi constru\'{i}do
explicitamente comparando formatos. Isso \'{e} mais correto: um par $(i,p,l)$ pode
existir no dicion\'{a}rio $k$ sem ser compat\'{i}vel em formato.
O Big-$M = 10^6$ garante que $w$ s\'{o} \'{e} acionado quando a f\'{a}brica esgota
toda capacidade f\'{i}sica.
\end{mdframed}

% ── 9.7 ─────────────────────────────────────────────────────────────────────
\subsection*{9.7 Restri\c{c}\~{o}es}

\begin{mdframed}[style=codigo]
\begin{lstlisting}
# R1: limite de armazenagem (espaco compartilhado entre SKUs)
@constraint(modelo, R1[p in P, l in L_p[p], t in T],
    sum(e[i,p,l,t] for i in I) <= e_max[(p,l)])

# R2: balanco de estoque (e[t=0] = 0 via operador ternario)
@constraint(modelo, R2[i in I, p in P, l in L_p[p], t in T],
    e[i,p,l,t] == (t == 1 ? 0.0 : e[i,p,l,t-1]) +
                  x[i,p,l,t] + g[i,p,l,t] - s[i,p,l,t] - a[i,p,l,t])

# R3: atendimento a demanda com fator de escalonamento
@constraint(modelo, R3[i in I, p in P, l in L_p[p], t in T],
    a[i,p,l,t] + w[i,p,l,t] == eta[(i,p,l,t)] * fator_demanda)

# R4: capacidade produtiva (taxa efetiva = rate * OEE; planta quebrada => 0h)
@constraint(modelo, R4[p in P, l in L_p[p], t in T],
    sum(x[i,p,l,t] / (r_taxa[(p,l,t)] * fator_oee)
        for i in I if (i,p,l) in C && r_taxa[(p,l,t)] > 0) <=
    (p == fabrica_quebrada ? 0.0 : 24.0))

# R5: caminhoes de saida por planta por dia
@constraint(modelo, R5[p in P, t in T],
    sum(f[i,p,d,t] / 200_000
        for i in I, d in P if d != p) <= limite_caminhoes)

# R6: doca de ENTRADA (frete chegando -> g redistribuindo pelas linhas)
@constraint(modelo, R6[i in I, p in P, t in T],
    sum(g[i,p,l,t] for l in L_p[p]) ==
    sum(f[i,o,p,t] for o in P if o != p))

# R7: doca de SAIDA (s coletando das linhas -> frete saindo)
@constraint(modelo, R7[i in I, p in P, t in T],
    sum(s[i,p,l,t] for l in L_p[p]) ==
    sum(f[i,p,d,t] for d in P if d != p))

# R8: capabilidade de formato (usa conjunto C, nao haskey)
@constraint(modelo, R8[i in I, p in P, l in L_p[p], t in T;
                        (i,p,l) not-in C],
    x[i,p,l,t] == 0)
\end{lstlisting}
\end{mdframed}

\begin{mdframed}[style=destaque]
\textbf{\color{ballblue}Notas sobre as restri\c{c}\~{o}es:}
\begin{itemize}[leftmargin=1.5em, itemsep=3pt]
  \item \textbf{R2 com \texttt{t == 1}:} o operador tern\'{a}rio implementa $e^0 = 0$
    sem separar em dois blocos de constraint.
  \item \textbf{R4 com \texttt{(i,p,l) in C}:} a vers\~{a}o anterior usava
    \texttt{haskey(k,(i,p,l))}. O c\'{o}digo correto usa o conjunto $\mathcal{C}$
    expl\'{i}cito baseado em formato.
  \item \textbf{R4 planta quebrada:} \texttt{p == fabrica\_quebrada ? 0.0 : 24.0}
    zera o limite de horas da planta, desligando-a completamente.
  \item \textbf{R4 efici\^{e}ncia (OEE):} o divisor \texttt{r\_taxa[(p,l,t)] * fator\_oee}
    aplica a efici\^{e}ncia operacional $\rho$. Com $\rho < 1$, a taxa efetiva cai e cada
    lata consome mais horas-m\'{a}quina, apertando R4.
  \item \textbf{R6 = entrada, R7 = sa\'{i}da:} os r\'{o}tulos est\~{a}o alinhados
    com o c\'{o}digo Julia. R6 redistribui o frete recebido ($f \to g$);
    R7 coleta o frete a expedir ($s \to f$).
  \item \textbf{R8 com $\mathcal{C}$:} \texttt{(i,p,l) \ensuremath{\notin} C}
    filtra as vari\'{a}veis que devem ser for\c{c}adas a zero por incompatibilidade de formato.
\end{itemize}
\end{mdframed}

% ============================================================
\newpage
\section{Painel de Controle e An\'{a}lise de Sensibilidade}

\subsection*{10.1 O que \'{e} An\'{a}lise de Sensibilidade em Pesquisa Operacional?}

\begin{mdframed}[style=destaque]
Um modelo de otimiza\c{c}\~{a}o resolve um problema para um conjunto \textbf{fixo} de par\^{a}metros.
Mas no mundo real, os par\^{a}metros mudam: a demanda cresce, equipamentos quebram,
caminhões entram em greve. A \textbf{An\'{a}lise de Sensibilidade} responde:

\begin{center}
  \textit{``At\'{e} quanto posso apertar um par\^{a}metro antes de a solu\c{c}\~{a}o \'{o}tima mudar drasticamente?''}
\end{center}

Em vez de refazer o modelo do zero, variamos os par\^{a}metros de controle e observamos
como o custo \'{o}timo, a taxa de atendimento e o uso log\'{i}stico respondem.
Essa \'{e} a base do \textbf{Painel de Controle} implementado no c\'{o}digo.
\end{mdframed}

\subsection*{10.2 Os quatro eixos de perturba\c{c}\~{a}o}

\begin{center}
\renewcommand{\arraystretch}{1.5}
\begin{tabular}{>{\bfseries\color{ballblue}}p{3.5cm} p{5cm} p{5cm}}
  \toprule
  Par\^{a}metro & O que representa & Como perturbar \\
  \midrule
  \texttt{fator\_demanda} &
    Multiplica toda a demanda $\eta_{i,p,l,t}$ por um fator. \newline
    $1.0$ = normal, $2.0$ = dobra a demanda. &
    Simula choques de mercado ou promo\c{c}\~{o}es. Combinado com \texttt{ocupacao\_100},
    equivale a testar ocupa\c{c}\~{o}es al\'{e}m dos 150\% dos dados. \\[4pt]
  \texttt{fabrica\_quebrada} &
    Reduz o limite de horas da planta escolhida de 24\,h para 0\,h em R4. &
    Simula parada de manuten\c{c}\~{a}o, greve ou desastre. Revela quais plantas
    s\~{a}o cr\'{i}ticas para a malha. \\[4pt]
  \texttt{limite\_caminhoes} &
    Substitui o valor 80 na restri\c{c}\~{a}o R5. &
    Simula greve de transportadores ou escassez de ve\'{i}culos.
    Revela o quanto a rede depende do frete interplanta. \\[4pt]
  \texttt{fator\_oee} &
    Efici\^{e}ncia operacional $\rho$ que multiplica a taxa nominal em R4
    (taxa efetiva $= r_{p,l,t}\cdot\rho$). $1.0$ = capacidade nominal. &
    Simula desgaste, setups e microparadas das linhas. Reduzir $\rho$ encurta
    a capacidade produtiva sem desligar a planta. \\
  \bottomrule
\end{tabular}
\end{center}

\subsection*{10.3 Estrat\'{e}gia de experimenta\c{c}\~{a}o sugerida}

O script roda os 21 cen\'{a}rios (\texttt{ocupacao\_50} a \texttt{ocupacao\_150}) com a
configura\c{c}\~{a}o fixa do painel. Para uma an\'{a}lise completa, sugere-se rodar em seq\^{u}encia:

\begin{center}
\renewcommand{\arraystretch}{1.4}
\begin{tabular}{c c c c p{5.5cm}}
  \toprule
  Rodada & \texttt{fator} & \texttt{fabrica} & \texttt{limite} & Objetivo \\
  \midrule
  1 & 1.0 & Nenhuma & 80 & Baseline --- comportamento normal da malha \\
  2 & 2.0 & Nenhuma & 80 & Choque de demanda --- onde aparece $w$? \\
  3 & 1.0 & Extrema & 80 & Planta cr\'{i}tica --- impacto no custo? \\
  4 & 1.0 & Nenhuma & 10 & Colapso log\'{i}stico --- a rede suporta? \\
  5 & 2.0 & Extrema & 10 & Cen\'{a}rio catastrófico --- at\'{e} onde vai $w$? \\
  \bottomrule
\end{tabular}
\end{center}

\subsection*{10.4 Lendo o output}

Cada rodada exporta um CSV nomeado automaticamente com a configura\c{c}\~{a}o usada,
contendo uma linha por cen\'{a}rio de ocupa\c{c}\~{a}o. Colunas-chave:

\begin{itemize}[leftmargin=1.5em, itemsep=3pt]
  \item \textbf{custo\_producao, custo\_estoque, custo\_logistica}: decomposi\c{c}\~{a}o do custo real
    --- identifica qual componente cresce mais com a ocupa\c{c}\~{a}o.
  \item \textbf{custo\_penalidade}: valor total pago em Big-$M$.
    Zero = 100\% de atendimento. Positivo = venda perdida.
  \item \textbf{taxa\_atendimento}: percentual da demanda efetivamente entregue ($a / \eta$).
  \item \textbf{caminhoes}: total de caminh\~{o}es despachados. Se atingir
    $80 \times 10 \times 15 = 12.000$, R5 est\'{a} saturada.
\end{itemize}

\subsection*{10.5 An\'{a}lise dual (pre\c{c}os-sombra)}

Al\'{e}m do resumo consolidado, cada cen\'{a}rio exporta o detalhamento primal
(\texttt{detalhe\_ocupacao\_<N>.csv}) e os \textbf{pre\c{c}os-sombra} das restri\c{c}\~{o}es
ativas --- a leitura dual do problema, que quantifica o valor marginal de relaxar cada
recurso. Como o modelo \'{e} um PL cont\'{i}nuo, os duais s\~{a}o obtidos diretamente do
solver (\texttt{shadow\_price} do JuMP) quando existe solu\c{c}\~{a}o dual vi\'{a}vel.

\begin{center}
\renewcommand{\arraystretch}{1.4}
\small
\begin{tabular}{>{\bfseries\color{ballblue}}c p{3.2cm} p{7.5cm}}
  \toprule
  Dual & Restri\c{c}\~{a}o & Interpreta\c{c}\~{a}o econ\^{o}mica \\
  \midrule
  $\mu_{p,l,t}$    & R4 (horas) &
    Varia\c{c}\~{a}o do custo \'{o}timo por $+1$ hora-m\'{a}quina na linha $l$.
    Exportado em \texttt{duals\_capacidade\_ocupacao\_<N>.csv} (\texttt{mu\_hora\_maquina}). \\[4pt]
  $\gamma_{p,l,t}$ & R1 (armazenagem) &
    Valor marginal de $+1$ unidade de capacidade de armazenagem
    (\texttt{gamma\_armazenagem} no mesmo arquivo). \'{E} $\approx 0$ quando o armazém n\~{a}o \'{e} gargalo. \\[4pt]
  $\pi_{i,p,l,t}$  & R3 (demanda) &
    Custo marginal de atender $+1$ lata de demanda. Exportado em
    \texttt{duals\_demanda\_ocupacao\_<N>.csv} (\texttt{pi\_demanda}), com a flag
    \texttt{saturado} indicando onde h\'{a} venda perdida ($w > 0$). \\
  \bottomrule
\end{tabular}
\end{center}

\begin{mdframed}[style=destaque]
\textbf{\color{ballblue}Como interpretar.} Onde a demanda \'{e} atendida via produ\c{c}\~{a}o
local com capacidade sobrando, $\pi$ reflete o custo unit\'{a}rio de produ\c{c}\~{a}o/estoque.
Onde R4 satura, $\mu$ torna-se relevante --- uma hora-m\'{a}quina extra ``vale''
a economia que ela gera. Nos pontos de \textbf{venda perdida}, $\pi$ aproxima-se do
Big-$M$ ($10^6$): a \'{u}ltima lata s\'{o} poderia ser atendida pagando a penalidade,
sinalizando que o gargalo \'{e} f\'{i}sico, n\~{a}o econ\^{o}mico.
\end{mdframed}

{\color{accentgray}\rule{\linewidth}{0.5pt}}
\begin{center}
  \small\color{darkgray}
  O c\'{o}digo completo est\'{a} no arquivo \texttt{modelagem.jl}. Ajuste os caminhos
  \texttt{input\_dir}/\texttt{output\_dir} no struct \texttt{Config} e as quatro vari\'{a}veis
  do Painel de Controle (\texttt{fabrica\_quebrada}, \texttt{limite\_caminhoes},
  \texttt{fator\_demanda}, \texttt{fator\_oee}) antes de rodar.
\end{center}

\end{document}